\section{Project Overview}

\subsection{Project Purpose, Scope, and Objectives}

Many patients, especially those with chronic diseases, take
multiple medications prescribed by different doctors.
These medications may interact with each other or contain
the same active ingredients, leading to potential health
risks. Different brand names and unclear prescriptions
make it difficult for patients to identify these risks
and manage their medication schedules.

The Medication Safety Assistant aims to help patients
identify potential medication risks and manage their
prescribed medication schedules. 
Users can upload prescription images, type natural-language
requests, use voice input, or enter medication details
through manual forms.
The system checks for drug interactions and duplicate
active ingredients, provides safety warnings, and supports
healthcare professional (HCP) review when needed.
It also offers medication reminders and adherence tracking
to help patients follow their confirmed schedules.


The project objectives are:

\begin{itemize}
    \item Support medication input through OCR, natural-language
    text, voice, and manual forms, with user confirmation
    before assessment.

    \item Identify user intent and extract relevant medication
    information from natural-language text.

    \item Identify Vietnamese medications and active ingredients
    using medicine reference data, and assess drug-drug
    interactions using DDInter 2.0.

    \item Detect duplicate active ingredients and provide
    clear risk levels, safety warnings, and data limitations.

    \item Support healthcare professional (HCP) review of
    high-risk or uncertain cases through a simulated workflow.

    \item Provide medication reminders based on confirmed
    prescription schedules, with safety checks and
    adherence tracking.

    \item Deliver a working Web prototype within five sprints.
\end{itemize}

The initial release excludes ADR reporting, drug-food
interaction checks, advanced medication timing analysis,
pill recognition, handwritten prescription support,
and integration with real healthcare systems.

ADR reporting may be added later if time allows,
after the core features are completed and tested.
The system does not independently prescribe, stop,
or change medications.

\subsection{Assumptions and Constraints}

\begin{longtable}{|L{0.18\textwidth}|L{0.35\textwidth}|L{0.35\textwidth}|}
\hline
\textbf{Type} & \textbf{Description} & \textbf{Effect on the project}\\ \hline

Assumption
& Users can review and correct medication information from OCR,
NLP, Voice, and Manual Input.
& Medication information must be confirmed before safety assessment.
\\ \hline

Assumption
& A limited Vietnamese medication dataset and
DDInter 2.0 data can be used for academic testing,
subject to their permitted terms.
& The team must check data access, reuse rights,
accuracy, and coverage before use.
\\ \hline

Assumption
& Healthcare professional (HCP) review can be simulated
for the prototype.
& The system will demonstrate the review process without
real clinical approval.
\\ \hline

Schedule constraint
& The project has three members and five two-week sprints.
& Core features must be prioritized to meet the deadline.
\\ \hline

Resource constraint
& Development uses student laptops, available datasets,
and free or approved tools.
& The prototype must work within limited resources.
\\ \hline

Technical constraint
& OCR and Voice may produce incorrect results.
Handwritten prescriptions are not supported.
& Users must be able to correct errors or enter medication
information manually.
\\ \hline

Data constraint
& Medication and interaction data may be incomplete.
& Missing information must be reported as unresolved
or unassessed, not considered safe.
\\ \hline

Safety constraint
& The system cannot independently prescribe, stop,
or change medications.
& Safety warnings must use supported data and recommend
professional advice when needed.
\\ \hline

Reminder constraint
& Users may confirm schedules from existing prescriptions.
Users may adjust reminder times within
confirmed prescription or HCP instructions.
Changes that conflict with or cannot be
checked against these instructions require
HCP review.
& Severe interactions must block affected new reminders
until review, without automatically stopping existing
or unrelated reminders.
\\ \hline

Privacy constraint
& Medication records, prescription images, and voice
data may contain sensitive information.
& The prototype must limit data collection, protect access,
and use synthetic or non-identifying test data.
\\ \hline

\end{longtable}


\subsection{Project Deliverables}

\begin{longtable}{|L{0.24\textwidth}|L{0.41\textwidth}|L{0.13\textwidth}|L{0.12\textwidth}|}
\hline
\textbf{Deliverable} & \textbf{Description / acceptance condition} & \textbf{Owner} & 
\textbf{Target}\\ \hline

PA1 planning package & Initial Project Plan, Vision Document, 3-6 UI screens, two weekly reports, and accessible tracking evidence. & Whole team & Sprint 1\\ \hline

Requirements Baseline
& Reviewed Vision Document, prioritized backlog,
use cases, and non-functional requirements.
& Long & Sprint 2\\ \hline

Architecture and UX baseline & System context, main workflow, data/interface decisions, and tested UI prototype covering OCR, manual, typed, and voice input. & Trang / Thành & Sprint 2\\ \hline

Medication Intake Proof of Concept
& OCR, NLP, Voice, and Manual Input with user confirmation
and a common medication data format.
& Whole team & Sprint 3\\ \hline

Medication Safety Proof of Concept
& Prepared Vietnamese medication data, reviewed ingredient
mappings to DDInter 2.0, and tested drug-drug interaction
and duplicate-ingredient checks with clear risk and
incomplete-data results.
& Long & Sprint 3\\ \hline

Integrated test build
& End-to-end Web application covering medication
input, identification, safety checks, warnings,
and simulated HCP review.
& Whole team & Sprint 4\\ \hline

Final prototype release
& Web prototype with confirmed-schedule reminders,
safety gate, adherence tracking, test results,
documentation, and demo materials.
& Whole team & Sprint 5\\ \hline

Project evidence repository & Versioned source, decisions, test evidence, sprint records, risk register, meeting notes, and weekly reports accessible to teaching staff. & Trang & Weekly\\ \hline
\end{longtable}

